\subsubsection{Green\textquotesingle s Functions for Boundary-Value Problems}\label{subsec:green-functions-boundary-problems}

We shall consider the boundary-value problem \eqref{eq:general-homogeneous-boundary-value-problem}, and we make the assumptions of the last section.

\begin{Definition}
\label{def:green-function-boundary-problem}
A \textbf{Green\textquotesingle s function} $G(\bm{x},\bm{y})$ for \eqref{eq:general-homogeneous-boundary-value-problem} is a solution of the problem
\begin{equation}\label{eq:green-function-boundary-value-definition}
\begin{aligned}
L(\bm{x},D_{\bm{x}})G(\bm{x},\bm{y})&=\delta(\bm{x}-\bm{y}),\quad\bm{x},\bm{y}\in\Omega,\\
B_j(\bm{x},D_{\bm{x}})G(\bm{x},\bm{y})&=0,\quad\bm{x}\in\partial\Omega,\ \bm{y}\in\Omega,\ j=1,\ldots,p.
\end{aligned}
\end{equation}

\end{Definition}

The first equation in \eqref{eq:green-function-boundary-value-definition} is to be interpreted in the sense of distributions. Of course, giving a meaning to the boundary conditions requires more smoothness of $G$ than that it be a distribution. For elliptic boundary-value problems, however, it turns out that $G$ is smooth as long as $\bm{x}\ne\bm{y}$, and hence the interpretation of the boundary conditions poses no problems. Clearly, the concept of a Green\textquotesingle s function generalizes that of a fundamental solution. If $L$ has constant coefficients, it is in fact often advantageous to think of the Green\textquotesingle s function as a perturbation of the fundamental solution. Namely, if $\mathcal{G}(\bm{x}-\bm{y})$ is the fundamental solution, we set $G(\bm{x},\bm{y})=\mathcal{G}(\bm{x}-\bm{y})+g(\bm{x},\bm{y})$ where $g$ satisfies
\begin{equation}\label{eq:green-fundamental-solution-smooth-perturbation}
L(\bm{x},D_{\bm{x}})g(\bm{x},\bm{y})=0,\quad\bm{x},\bm{y}\in\Omega,
\end{equation}
and for $j=1,\ldots,p$
\begin{equation}\label{eq:green-perturbation-boundary-data}
B_j(\bm{x},D_{\bm{x}})g(\bm{x},\bm{y})=-B_j(\bm{x},D_{\bm{x}})\mathcal{G}(\bm{x}-\bm{y}),\quad\bm{x}\in\partial\Omega,\ \bm{y}\in\Omega.
\end{equation}
If $\mathcal{G}$ is smooth for $\bm{x}\ne\bm{y}$, then the right-hand side of \eqref{eq:green-perturbation-boundary-data} is smooth for every $\bm{y}\in\Omega$. For elliptic problems, we shall see in Chapter~\ref{chap:linear-elliptic-equations} that this implies that $g$ is also smooth. In the interior of $\Omega$, the fundamental solution in a sense ``dominates'' the Green\textquotesingle s function by contributing the most singular part.

It is of course of fundamental importance to identify classes of boundary-value problems for which Green\textquotesingle s functions exist and are unique. At present, we do not have the techniques available which are required to address this question, but we shall address the issue of existence for elliptic equations in Chapter~\ref{chap:linear-elliptic-equations}.

If a Green\textquotesingle s function exists, then a formal solution of \eqref{eq:general-homogeneous-boundary-value-problem} is
\begin{equation}\label{eq:green-homogeneous-boundary-solution-representation}
u(\bm{x})=\int_\Omega G(\bm{x},\bm{y})f(\bm{y})\,\mathrm{d}\bm{y};
\end{equation}
in fact, if $f\in\mathcal{D}(\Omega)$, then \eqref{eq:green-homogeneous-boundary-solution-representation} gives a solution of \eqref{eq:general-homogeneous-boundary-value-problem} under fairly minimal assumptions on $G$. It suffices, for example, if $G(\cdot,\bm{y})$ as an element of $\mathcal{D}^{\prime}(\Omega)$ depends continuously on $\bm{y}$ and $G$ is smooth for $\bm{x}\ne\bm{y}$. In particular, if the boundary-value problem \eqref{eq:general-homogeneous-boundary-value-problem} is uniquely solvable, \eqref{eq:green-homogeneous-boundary-solution-representation} leads to the identity
\begin{equation}\label{eq:green-test-function-representation-identity}
u(\bm{x})=\int_\Omega G(\bm{x},\bm{y})L(\bm{y},D_{\bm{y}})u(\bm{y})\,\mathrm{d}\bm{y}
\end{equation}
for all $u\in\mathcal{D}(\Omega)$.

We shall now assume that $G$ has sufficient regularity to establish \eqref{eq:green-homogeneous-boundary-solution-representation} not only for $u\in\mathcal{D}(\Omega)$, but for every $u\in X_B$. Using \eqref{eq:general-boundary-operator-green-identity}, we conclude that
\begin{equation}\label{eq:green-adjoint-kernel-boundary-derivation}
\begin{aligned}
u(\bm{x})&=\int_\Omega u(\bm{y})L^*(\bm{y},D_{\bm{y}})G(\bm{x},\bm{y})\,\mathrm{d}\bm{y}\\
&\quad+\int_{\partial\Omega}\sum_{j=1}^p S_j(\bm{y},D_{\bm{y}})u(\bm{y})C_j(\bm{y},D_{\bm{y}})G(\bm{x},\bm{y})\,\mathrm{d}S_{\bm{y}}.
\end{aligned}
\end{equation}
If this holds for arbitrary $u\in X_B$, we find that, for every $\bm{x}\in\Omega$, we must have
\begin{equation}\label{eq:green-kernel-adjoint-variable-problem}
\begin{aligned}
L^*(\bm{y},D_{\bm{y}})G(\bm{x},\bm{y})&=\delta(\bm{x}-\bm{y}),\quad\bm{x},\bm{y}\in\Omega,\\
C_j(\bm{y},D_{\bm{y}})G(\bm{x},\bm{y})&=0,\quad\bm{y}\in\partial\Omega,\ \bm{x}\in\Omega,\ j=1,\ldots,p.
\end{aligned}
\end{equation}
That is, $G$, regarded as a function of $\bm{y}$ for fixed $\bm{x}$, satisfies the adjoint boundary-value problem.

Using \eqref{eq:green-kernel-adjoint-variable-problem} and setting $v(\bm{y})=G(\bm{x},\bm{y})$ in \eqref{eq:general-boundary-operator-green-identity}, we find
\begin{equation}\label{eq:green-inhomogeneous-boundary-representation-identity}
\begin{aligned}
u(\bm{x})&=\int_\Omega G(\bm{x},\bm{y})L(\bm{y},D_{\bm{y}})u(\bm{y})\,\mathrm{d}\bm{y}\\
&\quad+\int_{\partial\Omega}\sum_{j=1}^p T_j(\bm{y},D_{\bm{y}})G(\bm{x},\bm{y})B_j(\bm{y},D_{\bm{y}})u(\bm{y})\,\mathrm{d}S_{\bm{y}}.
\end{aligned}
\end{equation}
Thus, if the inhomogeneous boundary-value problem
\begin{equation}\label{eq:general-inhomogeneous-boundary-value-problem}
\begin{aligned}
L(\bm{x},D)u&=f(\bm{x}),\quad\bm{x}\in\Omega,\\
B_j(\bm{x},D)u&=\phi_j(\bm{x}),\quad\bm{x}\in\partial\Omega,\quad j=1,\ldots,p
\end{aligned}
\end{equation}
has a solution, then the solution is represented by
\begin{equation}\label{eq:green-general-inhomogeneous-solution-formula}
\begin{aligned}
u(\bm{x})&=\int_\Omega G(\bm{x},\bm{y})f(\bm{y})\,\mathrm{d}\bm{y}\\
&\quad+\int_{\partial\Omega}\sum_{j=1}^p T_j(\bm{y},D_{\bm{y}})G(\bm{x},\bm{y})\phi_j(\bm{y})\,\mathrm{d}S_{\bm{y}}.
\end{aligned}
\end{equation}
As a caution, we note that in justifying the integration by parts which leads to \eqref{eq:green-inhomogeneous-boundary-representation-identity}, it is important that $\bm{x}\in\Omega$ so that $G(\bm{x},\bm{y})$ is smooth for $\bm{y}\in\partial\Omega$. In general, \eqref{eq:green-inhomogeneous-boundary-representation-identity} does not represent $u(\bm{x})$ for $\bm{x}\in\partial\Omega$.

For some simple equations in simple domains, Green\textquotesingle s functions can be given explicitly. As an example, we consider Laplace\textquotesingle s equation on the ball $B_R$ of radius $R$ with Dirichlet boundary conditions. In this case, the Green\textquotesingle s function can be constructed by what is known as the method of images. The fundamental solution $\mathcal{G}(|\bm{x}-\bm{y}|)$ can be thought of as the potential of a point charge located at the point $\bm{y}$. The idea is now to put a second point charge at the reflected point $\overline{\bm{y}}=R^2\bm{y}/|\bm{y}|^2$ in such a way that the potentials of the two charges cancel each other on the sphere $|\bm{x}|=R$. This leads to the Green\textquotesingle s function
\begin{equation}\label{eq:ball-dirichlet-green-image-formula}
\begin{aligned}
G(\bm{x},\bm{y})&=\mathcal{G}(|\bm{x}-\bm{y}|)-\mathcal{G}\left(\dfrac{|\bm{y}|}R|\bm{x}-\overline{\bm{y}}|\right)\\
&=\mathcal{G}\left(\sqrt{|\bm{x}|^2+|\bm{y}|^2-2\bm{x}\cdot\bm{y}}\right)\\
&\quad-\mathcal{G}\left(\sqrt{(|\bm{x}||\bm{y}|/R)^2+R^2-2\bm{x}\cdot\bm{y}}\right).
\end{aligned}
\end{equation}
If $\bm{y}=0$, we set $G(\bm{x},\bm{y})=\mathcal{G}(|\bm{x}|)-\mathcal{G}(R)$. If $|\bm{x}|=R$, then $G(\bm{x},\bm{y})=0$; moreover, we compute
\begin{equation}\label{eq:ball-green-image-laplacian}
\Delta_{\bm{x}}G(\bm{x},\bm{y})=\delta(\bm{x}-\bm{y})-\dfrac{|\bm{y}|^2}{R^2}\delta(\bm{x}-\overline{\bm{y}}),
\end{equation}
which agrees with $\delta(\bm{x}-\bm{y})$ if $\bm{x}$ and $\bm{y}$ are restricted to $B_R$. Hence $G$ is indeed a Green\textquotesingle s function for the Dirichlet problem. We see that $G$ is symmetric in its arguments, reflecting the self-adjointness of the Dirichlet problem.

The solution of the Dirichlet problem
\begin{equation}\label{eq:ball-dirichlet-poisson-boundary-problem}
\Delta u=f(\bm{x}),\ \bm{x}\in B_R,\qquad u=\phi(\bm{x}),\ \bm{x}\in\partial B_R
\end{equation}
is represented by \eqref{eq:green-general-inhomogeneous-solution-formula}:
\begin{equation}\label{eq:ball-dirichlet-green-solution-representation}
u(\bm{x})=\int_{B_R}G(\bm{x},\bm{y})f(\bm{y})\,\mathrm{d}\bm{y}+\int_{\partial B_R}\dfrac{\partial}{\partial n_{\bm{y}}}G(\bm{x},\bm{y})\phi(\bm{y})\,\mathrm{d}S_{\bm{y}}.
\end{equation}
Moreover, a direct calculation shows that
\begin{equation}\label{eq:ball-green-boundary-normal-derivative}
\dfrac{\partial}{\partial n_{\bm{y}}}G(\bm{x},\bm{y})=\dfrac{R^2-|\bm{x}|^2}{\Omega_m R|\bm{x}-\bm{y}|^m}
\end{equation}
(cf. Problem~\ref{ex:verify-ball-green-normal-derivative}). This leads to Poisson\textquotesingle s formula, which we have already encountered in Section~\ref{sec:dirichlet-existence-proof}.
